\scanpage{138}
\[
 z_1=u(e_1),\ z_2=u(e_2)\in H\subset V
\]
alors la structure de Torelli $u\circ g$ correspond à $z'_1,z'_2\in\Pi$ donnée par
\[
 \begin{aligned}
 z'_1&=ug(e_1)=u(ae_1+ce_2)=az_1+cz_2\\
 z'_2&=ug(e_2)=u(be_1+de_2)=bz_1+dz_2,
 \end{aligned}
\]
qui correspond donc à un élément
\[
 z'=\unsure{\alpha'|x'}+iy'\in\mathscr D'
\]
donné par
\[
 z'\deleted{\underbrace{\unsure{=x'+iy'|=-x'+iy'}}_{\unsure{\tilde z'|z'}\; z'=-\bar z'}}
 =-\frac{z'_2}{z'_1}=-\frac{bz_1+dz_2}{az_1+cz_2}
 =-\frac{d\frac{z_2}{z_1}+b}{c\frac{z_2}{z_1}+a}
 =\frac{d\unsure{z|\bar z}\unsure{-|+}b}{-c\unsure{z|\bar z}+a}
 \tag*{\deleted{\unsure{(24)}} (25)}
\]
% slanted
\margin{\deleted{d'où $-z'=\frac{-dz+b}{-cz+a}$ \quad (25) $z'=\frac{dz-b}{-cz+a}$}}
où \deleted{$\unsure{\tilde z=-x+iy=-\bar z}$.} Si on \deleted{\unsure{reg}} préfère expliciter \deleted{l'}l'action à gauche $u\longmapsto u\circ g^{-1}$ de $\mathrm{Sl}(2,\mathbb Z)$ sur l'ens. des structures de Torelli mod iso, on note que
\[
 g^{-1}=\begin{pmatrix}d&-b\\-c&a\end{pmatrix}\qquad \text{\deleted{d'où,}}
\]
d'où maintenant, pour $z'=x'+iy'$ correspondant à $u\circ g^{-1}$
\[
 \deleted{\tilde z'=\bar z'=\frac{a\tilde z-b}{-c\tilde z+d}
 =\frac{-a\bar z-b}{-c\bar z-d}
 =\frac{\overline{a\bar z+b}}{\overline{c\bar z+d}}} \tag*{\deleted{\unsure{(25bis)}}}
\]
\[
 z'=\frac{az+b}{cz+d} \tag*{(25bis)}
\]
\uline{victoire}, c'est la formule habituelle, mais on fera attention que le nb\textsuperscript{re} complexe $z\in\mathscr D$ attaché à la structure de Torelli définie par \deleted{$z_1,z_2\in\mathfrak t_E$}\inserted{les deux} éléments $z_1,z_2$ de la droite complexe \deleted{$\unsure{\mathfrak t_{\mathbb R}=V}$} $\mathfrak t_E$ est donné, non
% archive sheet number: 137
